\section{Introdução}
\label{sec:introducao}

\subsection{Contexto: reconhecimento facial como infraestrutura social}
\label{sec:contexto}

Sistemas de reconhecimento facial deixaram, na última década, de ser uma tecnologia restrita a laboratórios de pesquisa para integrar a infraestrutura cotidiana de identificação humana. Aplicam-se em desbloqueio de dispositivos pessoais, autenticação bancária e identidade digital, controle de acesso a ambientes corporativos, controle de fronteiras em aeroportos e identificação policial em vias públicas. Em escala industrial, o programa \emph{Face Recognition Vendor Test} do National Institute of Standards and Technology (NIST) conduziu, em 2019, sob a coordenação de \textcite{grother2019face}, a maior auditoria pública já realizada em biometria facial --- cento e oitenta e nove algoritmos comerciais avaliados sobre dezoito milhões de imagens de oito milhões e meio de pessoas.

Do lado científico, a maturação técnica também se acelerou. Modelos modernos combinam grandes datasets faciais~\autocite{karkkainen2021fairface,wang2019racial,robinson2020face}, arquiteturas neurais profundas fundadas em ResNet, ConvNeXt e Vision Transformers, e funções de perda especializadas (ArcFace, CosFace), atingindo desempenho próximo ao humano em condições controladas. Em paralelo, \emph{vision-language models} como PaliGemma reformularam a classificação racial como \emph{Visual Question Answering}: o estado da arte atual em classificação racial multi-classe sobre o FairFace, o \textbf{FaceScanPaliGemma}~\autocite{aldahoul2024exploring}, atinge F1 macro de 75\,\% e acurácia agregada de 75,7\,\%.

Esta pesquisa parte de uma constatação simples e bem documentada: o reconhecimento facial funciona bem para a maioria das pessoas, mas não funciona igualmente bem para todas. A tecnologia amadureceu; a garantia de equidade, não. É sobre essa diferença que a dissertação se debruça.

\subsection{Problema: disparidade demográfica residual}
\label{sec:problema}

O problema central que motiva a presente pesquisa enuncia-se de forma direta: \textbf{sistemas de reconhecimento facial apresentam disparidade demográfica significativa em acurácia entre grupos raciais, mesmo quando treinados sobre datasets explicitamente balanceados por design}.

\subsubsection{Disparidade documentada no estado da arte}
\label{sec:problema_disparidade}

Apesar do F1 macro de 75\,\% reportado pelo FaceScanPaliGemma~\autocite{aldahoul2024exploring}, o detalhamento por classe racial revela disparidade severa: aproximadamente 90\,\% para faces classificadas como \emph{Black}, 80\,\% para \emph{White}, 67\,\% para \emph{Southeast Asian} e apenas 60\,\% para \emph{Latinx/Hispanic}. A pior classe (Latinx) exibe F1 aproximadamente trinta pontos percentuais abaixo da melhor classe (Black).

A disparidade não é artefato de um único modelo. O baseline canônico ResNet-34 originalmente treinado por \textcite{karkkainen2021fairface} sobre o FairFace apresenta padrão análogo: 72\,\% de acurácia agregada, porém apenas 59,6\,\% em Hispanic contra 83,2\,\% em Black, conforme reportado independentemente por \textcite{lin2022fairgrape} em experimentos de pruning \emph{fairness-aware}. A hierarquia de dificuldade entre raças mostra-se \textbf{estável através de múltiplos métodos, arquiteturas e datasets}, o que sugere um problema estrutural em vez de algorítmico contingente.

\subsubsection{Balanceamento de dados não resolve}
\label{sec:problema_balanceamento}

A solução intuitiva de balancear o conjunto de treino por raça já foi exaustivamente testada e demonstrou-se insuficiente. O próprio FairFace foi construído com esse propósito~\autocite{karkkainen2021fairface}, distribuindo suas 108.501 imagens de forma aproximadamente equilibrada entre sete categorias raciais. Mesmo assim, a disparidade Latinx \emph{versus} White persiste.

\textcite{kolla2022impact} conduziram dezesseis experimentos com proporções variáveis de raças no treino, concluindo explicitamente que ``distribuição uniforme de raças nos dados de treinamento não garante face recognition livre de viés''. \textcite{wang2019racial}, ao introduzir o dataset \emph{Racial Faces in-the-Wild} (RFW), observaram que ``algumas raças permanecem intrinsecamente mais difíceis de reconhecer, mesmo com dados de treinamento balanceados''. Em escala industrial, o mesmo relatório NISTIR 8280 de \textcite{grother2019face} documentou diferenciais de dez a cem vezes na taxa de falso positivo entre grupos raciais em algoritmos comerciais.

\subsubsection{Diagnóstico: heterogeneidade fenotípica intra-categorial}
\label{sec:problema_heterogeneidade}

Trabalhos recentes apontam duas direções complementares para explicar essa disparidade residual. A primeira, empírica, sustenta que a categoria \emph{Latinx/Hispanic} do FairFace mascara heterogeneidade fenotípica interna substantiva. Evidência convergente vem de três disciplinas independentes: a antropologia biológica --- via o projeto PERLA (\emph{Project on Ethnicity and Race in Latin America}) coordenado por \textcite{telles2014pigmentocrac}, que documentou pigmentocracia em quatro países latino-americanos ---, a genética populacional --- com o estudo de \textcite{bryc2015genetic} sobre 162.721 indivíduos, evidenciando composição altamente variável de ancestralidade Native American, European e African em Latinos ---, e a sociologia identitária --- com a análise longitudinal do Pew Research Center~\autocite{lopez2017hispanic}, que documenta declínio da identidade Hispanic de 97\,\% para 50\,\% ao longo de quatro gerações nos Estados Unidos. A rotulagem monolítica no FairFace, portanto, agrupa em uma única categoria fenótipos que atravessam a totalidade do espectro Monk Skin Tone.

A segunda direção, também recente, sugere que o \emph{tom de pele} --- e não a categoria racial nominal --- é o \emph{driver} estrutural do gap. \textcite{pangelinan2023exploring} apresentam análise causal que atribui a disparidade em face recognition primariamente a diferenças de \emph{pixel information} (fração de face útil na imagem), colocando o balanceamento racial em segundo plano. \textcite{matias2026large} confirmam a maturidade do tom de pele como dimensão auditável ao publicar o \textbf{SkinToneNet}, primeiro classificador MST em larga escala.

\subsubsection{Implicações sociais e regulatórias}
\label{sec:problema_sociais}

A disparidade tem consequências concretas: quando um sistema é utilizado em catraca biométrica, autenticação financeira ou identificação policial, \textbf{a probabilidade de erro depende do rosto da pessoa}. Casos documentados de prisão errada por falso \emph{match} em sistemas comerciais de \emph{face recognition} já resultaram em litígios e moratórias regulatórias em múltiplas jurisdições. O \emph{European AI Act}, publicado em 2024, incluiu auditoria de \emph{fairness} como requisito formal para sistemas classificados como de alto risco~\autocite{lafargue2025fairness}.

A questão científica, portanto, deixou de ser ``existe viés?'' --- está documentada desde o estudo seminal \emph{Gender Shades}~\autocite{buolamwini2018gender}. A questão atual é: \textbf{qual é o mecanismo gerador desse viés residual mesmo em datasets balanceados, e como mitigá-lo de forma cientificamente defensável?}

\subsection{Objetivo e contribuições da pesquisa}
\label{sec:intro_objetivo}

Propõe-se desenvolver e avaliar um pipeline de classificação racial em imagens faciais que incorpora o tom de pele explícito (escala Monk Skin Tone, dez classes) como sinal auxiliar condicionante, por meio do mecanismo arquitetural \emph{Feature-wise Linear Modulation} (FiLM)~\autocite{perez2018film}, aplicado sobre \emph{backbone} ConvNeXt-T. A escolha da escala MST em detrimento da escala Fitzpatrick fundamenta-se no viés fototípico documentado desta última~\autocite{fitzpatrick1988validity} e na consolidação da MST como padrão moderno para auditoria de \emph{fairness}~\autocite{schumann2023consensus,matias2026large}. A escolha do mecanismo FiLM em detrimento de concatenação simples fundamenta-se em sua capacidade de modular \emph{features} intermediárias por escala e deslocamento, discutida em detalhe na Seção~\ref{sec:film} do Capítulo de Metodologia.

Três contribuições principais são pretendidas:

\begin{enumerate}[label=(\roman*)]
\item a primeira matriz pública de distribuição cruzada MST~$\times$~classes raciais sobre o FairFace \emph{validation set};
\item a primeira instância empírica documentada do uso de tom de pele como contexto arquitetural em race classification multi-classe;
\item uma decomposição quantitativa do erro Latinx em componente fenotípico (irredutível, derivado do \emph{overlap} MST intra-categoria) e componente algorítmico (mitigável, atacável por conditioning arquitetural), oferecendo diagnóstico estrutural no lugar de mera redução agregada de F1.
\end{enumerate}

A pesquisa é conduzida em conformidade com o \emph{European AI Act} (2024), que estabelece obrigatoriedade de auditoria de \emph{fairness} em sistemas biométricos classificados como de alto risco.

\subsection{Organização da dissertação}
\label{sec:intro_organizacao}

O restante deste documento está organizado como segue. O Capítulo~\ref{sec:rl} apresenta a revisão de literatura em seis frentes e identifica cinco lacunas científicas endereçadas por esta dissertação. O Capítulo~\ref{sec:objetivos} formaliza o objetivo geral, seis objetivos específicos, seis hipóteses testáveis e sete contribuições esperadas. O Capítulo~\ref{sec:metodologia} descreve o pipeline metodológico em seis etapas, as quatro configurações comparativas de \emph{conditioning}, os seis \emph{baselines} de mitigação, a triangulação de métricas fundamentada no Teorema da Impossibilidade~\autocite{kleinberg2017inherent} e o enquadramento ético do projeto no Art.~8\textsuperscript{o} da Resolução n\textsuperscript{o} 200/2021/CONSU da Unifesp. O Capítulo~\ref{sec:cronograma} apresenta o cronograma de execução, riscos identificados e próximos passos.
